\section{Five defects that were ours}
\label{app:ours}

These five fail the test the paper applies throughout, since a different group
writing correct code would not meet them. They appear here because they
circulated internally and the record of what got published matters more than a
tidy account.

\begin{enumerate}\itemsep3pt
\item \textbf{A leakage guard blind to its own namespace.} The invariant barring
a measurement target from defining the labels it gets measured against was
checked against a symbol list unable to contain the identifiers under
comparison. The guard passed on every input, including inputs violating it.

\item \textbf{An acceptance test over two disjoint identifier spaces.} It
compared gene symbols against database accessions. An empty intersection
follows by construction, so the assertion held on every run.

\item \textbf{A de-duplication that collapsed two population definitions.} A
per-patient de-duplication took the first row per patient from a frame carrying
\emph{two} pre-registered definitions of the diseased arm, selecting one by
file order. Measured, the same one in 224 of 224 combinations. Neither the
result table nor its sidecar recorded which. Five patients under the selected
definition held zero cells of the relevant compartment for a disease-calling
reason and entered the statistic as guaranteed abstentions, turning a
cohort-design fact into a positivity finding.

\item \textbf{A join that validated its keys and never its counts.} We assembled
the primary decomposition table by joining per-patient \emph{means} from one
job onto per-patient \emph{fractions} from another. The two jobs label cells at
different quality-control thresholds, so the populations differ by a median
factor of 1.26 and 1.35 on the two arms, and every term in
Equation~\ref{eq:kitagawa} multiplied a fraction over one set of cells by a
mean over another. The producing job's documentation asserts a loop mirroring
the other pipeline's filters and enumerates them. The threshold is the one
filter the loop misses. The replacement derives fractions, means, counts and
estimability from a single labelling call, then asserts its counts against the
committed counting function, which is the one line missing before.

\item \textbf{A provenance stamp blind to uncommitted code.} Every result must
carry a commit hash and a fixed seed. The dirtiness check guarding the hash
ignored untracked files, so a table produced by a script never committed
recorded a clean hash pointing at a commit lacking the script. In one case the
script entered the repository 83 seconds \emph{after} the table it wrote got
stamped clean against the previous commit. The check now counts untracked files
and the sidecar names them, tested against a repository holding an uncommitted
producing script --- the input that made the old check return clean.
\end{enumerate}

Items 3 and 4 are one failure twice over, a statistic computed over a
population nobody chose, and someone trying to prevent exactly that failure
wrote both. The de-duplication cites the violated invariant by name in its own
comment. Neither felt wrong, and both survived ordinary review.

\section{The bound, derived}
\label{app:bound}

\S\ref{sec:ceiling} rests on one inequality, so here it is in full rather than
by citation. \citet{chengliu2016} give the maximal point-biserial correlation
under several non-normal distributions including the uniform; the uniform case
is the one a rank statistic puts you in, and it is four lines.

Let $Y$ be binary with $\Pr(Y = 1) = p$, and let $X$ be the \emph{ranks} of the
continuous variable, so $X \sim \mathrm{Uniform}(0,1)$ in the limit. Then
$\mathrm{Var}(X) = 1/12$ and $\mathrm{Var}(Y) = p(1-p)$. The covariance is
maximised by perfect separation, with the $p$-fraction carrying $Y = 1$
occupying the top $p$ of the ranks, giving $\mathbb{E}[XY] = p\,(1 - p/2)$ and
\[
\mathrm{Cov}(X, Y) \;=\; p\left(1 - \tfrac{p}{2}\right) - \tfrac{1}{2}p
\;=\; \tfrac{1}{2}p(1-p),
\qquad
\rho_{\max} \;=\; \frac{\tfrac{1}{2}p(1-p)}
{\sqrt{\tfrac{1}{12}\,p(1-p)}} \;=\; \sqrt{3p(1-p)}.
\]
Two consequences the paper uses. Solving $\sqrt{3p(1-p)} = 0.20$ gives $p =
1.3516\%$: below that prevalence a tolerance of $0.20$ is unreachable and the
check cannot fire on any data. And $\rho_{\max} \to 0$ as $p \to 1$ just as it
does as $p \to 0$, so the check is blind at \emph{both} ends of its range.

We verified the identity numerically against perfect separation at $p = 0.005$,
$0.013516$, $0.05$, $0.10$, $0.30$ and $0.50$ with two million points each,
agreeing to better than $10^{-13}$ at every prevalence. The regression test in
the repository asserts it at four prevalences, including the $0.866$ ceiling at
$p = 0.5$ that surprises people.


\section{Three alternatives that look like escapes}
\label{app:alternatives}

\S\ref{sec:depth} states the short form for two of these. All three deserve
the full argument, because each is the first thing some reader reaches for.

\paragraph{Count splitting and data thinning.} \citet{neufeld2023} and the
data-thinning family split a cell's UMI counts into independent folds, so the
same genes can label a cell in one fold and be tested in the other. That
dissolves the leakage constraint as a \emph{statistical} matter and is the
right tool for the double-dipping problem it was built for. It does not rescue
us: splitting halves an already depth-limited count, and the label's dependence
on depth is a property of \emph{detection}, which survives the split and
worsens with it. A cell too shallow to show a marker in the full count is more
so in half of it. Our constraint is not that the label is statistically
dependent on the test data. It is that the label is a reading of how much data
there was.

\paragraph{Positive markers that are not targets.} The third escape, and the
first one a reader of \S\ref{sec:depth} reaches for, is to label the mature state
by positive markers held \emph{outside} the measured set --- AQP8, CEACAM7,
SLC26A3, KRT20 and the rest of our panel's exploratory tier. Our panel
forecloses it by construction: those genes are in the panel, and the leakage
invariant withholds the whole panel rather than the outcome gene alone. That is
a design choice, and a defensible one --- a label built from genes the same
disease process regulates re-imports the confound through a second door, and
tier A's own behaviour in \S\ref{sec:controls} is what that looks like --- but
it is \emph{ours}, not a property of the assay. A panel scoped to a single
outcome gene would not face it, and the honest statement of \S\ref{sec:depth}'s
premise is the conditional one: the move into absence is forced for any panel
where a state's positive markers double as the quantity of interest.

\paragraph{Reference-based annotators.} SingleR \citep{aran2019}
and CellTypist \citep{dominguezconde2022} score a cell against whole expression
profiles rather than a handful of binary marker calls, which does reduce the
fragility of absence-defined labels. The leakage invariant still applies ---
the target genes must be withheld from the reference --- and what remains of
the profile is measured through the same detection process. A softer label is
not a depth-independent one. We regard this as the most promising direction and
not as a solved problem, and we did not use it here.

\section{The tier collapse, and the grid the counts are over}
\label{app:collapse}

\paragraph{The algebra.} With $\mathrm{comp} = \Delta f\, m_N$ and
$\mathrm{intr} = f_N (m_T - m_N)$, the ratio between the arms equals
$(f_N/\Delta f)(m_T/m_N - 1)$. The gene enters only through $m_T/m_N$. For a
genuinely retained gene $m_T/m_N \approx 1$, the bracket goes to zero and tier D
separates cleanly, so the collapse is not algebraically forced in general. It is
forced once $m_T/m_N \approx 0$ for every tier. Measured per patient, the ratios
then agree to within 3\% across tiers, which is what the bracket ranging over
$[-1, -0.95]$ predicts --- the observed spread is the prediction landing, not a
precision the theory forbids.

\paragraph{Every row of \S\ref{sec:ceiling}, by arm.} The replication cohort's
per-arm table is 10 patients $\times$ 2 labelling axes $\times$ 4 resolutions
$\times$ 2 arms $=$ 160 rows. Of these, 102 carry a finite correlation and 58 do
not: 18 are the diseased arm at the rarest resolution, every one at prevalence
exactly $0$, and 40 are the coarsest resolution in both arms, every one at
prevalence exactly $1$. \textbf{No row lies in $0 < p < 1.3516\%$ or in
$98.65\% < p < 1$.} The ratio $|\rho|$-to-bound is defined on 22 of the 40 rows
at the rarest resolution (20 reference-arm, 2 diseased) and on all 40 at each of
the two middle resolutions; it is undefined at the coarsest, where the bound is
zero.

\paragraph{The grid the counts of \S\ref{sec:ceiling} are over.} One grid
throughout: patients $\times$ annotation resolutions $\times$ arms. Denominators
differ because it is filtered differently for different questions. 64 is one
resolution across both annotation axes on the primary cohort (GSE178341,
$n = 32$). The per-arm counts out of 160 and 40 are the replication cohort
(GSE132465, $n = 10$), which is what Figure~\ref{fig:ceiling} plots; the
per-subtype counts of Appendix~\ref{app:rare} are that cohort and GSE144735.

\section{The abundance-matched null, and what the simulation showed}
\label{app:floor}

We simulated both genes of \S\ref{sec:floor} \emph{purely compositional} --- no
intrinsic component anywhere, the same \emph{relative} patient-to-patient
variation on each --- so the pre-registered comparison starts as a coin flip.
With no assay noise the low-abundance gene scores below the high-abundance one in
51\% of 200 draws, as it should. Adding a shared additive noise floor and
changing nothing else takes the primary arm to \emph{every} draw at a floor of
0.5 units, and both arms to every draw at 2. The severity therefore depends on a
noise parameter nobody knows before the run, which is why no threshold on that
statistic could have been given a meaning in advance.

The scale-free fallback proposed at the time, a log-log slope, was tested and
carries the same confound; it is recorded in the configuration's
\texttt{withdrawn\_arms} so nobody re-proposes it.

The repair is to score each gene against an abundance-matched null: draw
off-panel genes at the outcome gene's own expression level by a rule fixed in
advance, and read the gene's percentile within that null rather than its raw
value or its rank against another gene.

\section{What the repairs recover, and from where}
\label{app:priorart}

Each repair this paper proposes has a literature, and none of it is ours.

\paragraph{Testing whether a label tracks a technical covariate.} The field's
tools are kBET \citep{buttner2019}, LISI \citep{korsunsky2019} and the scIB
suite \citep{luecken2022}, which score the covariate's local structure rather
than thresholding a correlation. Had we reached for one, \S\ref{sec:ceiling}
would not exist. What the bound adds is the reason a correlation threshold is the
wrong instrument, which is worth stating because a correlation threshold is what
gets written when nobody reaches for the standard tool. The rank-biserial
correlation $2\,\mathrm{AUC} - 1$ is the minimal fix inside the same idiom: it is
a monotone transform of the Mann--Whitney statistic, and under perfect separation
it equals 1 at every prevalence, which we verified numerically at
$p = 0.002$ through $0.5$ against the Spearman values that trace
$\sqrt{3p(1-p)}$ exactly.

\paragraph{Scoring a gene against an abundance-matched null.} Expression-matched
background sets are standard in gene-set analysis, and the competitive versus
self-contained distinction \citep{wu2012} is exactly the one \S\ref{sec:floor}
needed and did not make: our arm compared one gene against another gene rather
than against a background drawn at its own abundance.

\paragraph{Control genes as an estimator of unwanted variation.} Negative-control
methodology \citep{gagnonbartsch2012} treats control genes as an object to be
estimated from rather than a binary pass. Our tier system asks a small hand-picked
panel to answer yes or no, which is what makes a mis-specified tier
(\S\ref{sec:controls}) indistinguishable from a broken estimator.

\section{Extended notes}
\label{app:extended}

\paragraph{The unit of inference is the patient.} All bootstrapping resamples
patients, never cells --- the pseudoreplication point of \citet{squair2021}. An
earlier internal result violated this and reported a cell-pooled difference of
$+25.1\%$ where a single patient supplied 1{,}824 of 5{,}564 cells. Recomputed
per patient the difference is $+0.170$ with a 95\% interval of $[-0.011,
+0.342]$, containing zero. The direction survives independent 1:1
within-patient depth matching at $+0.182$, making it a direction with support
rather than an established effect.

\paragraph{Estimation runs per study, and the meta-analysis does not exist.}
Batch correction removes between-dataset variation, which is where the
compositional signal lives, so estimates run per study rather than pooled. Our
design named meta-analytic intervals as an output, and we now report that the
data cannot supply them. A meta-analysis needs two quotable (study, resolution)
cells and this project has one: the primary cohort's middle resolution on the
depth-matched read. Every other cell fails for a reason matching leaves
untouched. One coarser resolution is a two-bin split on about 90\% of patients,
the coarsest is degenerate by construction, the finest carries a median of one
usable cell, and the replication cohort's middle resolution stays confounded at
a median $|\rho| = 0.26$ in its reference arm. At $k = 1$ an engine combining
one study returns a number looking meta-analytic without being one. The cohort
holds one place to stand, which is the non-identifiability finding one level
up.

\paragraph{The two draw pools of \S\ref{sec:abstention} differ in more than one
way.} Neither matches the population the estimator meets in real data. There
the reference arm holds reference-tissue cells and the comparison arm does not,
whereas the second reading puts reference-tissue cells into \emph{both}. A
generator drawing each arm from its own tissue is a third option we did not
implement, because the imposed shift would then sit on top of a real
between-tissue difference and the truth would stop being known. The pools also
differ in target-gene dispersion, in how many distinct cells a replicate draws
from, and in whether the mature-cell distribution is unimodal or a mixture. The
sweep varies all three together and identifies none. A dispersion-only account
predicts discrimination rising with $\sqrt{n}$, and neither curve does.

\paragraph{Replicates are not independent, and more of them do not fix it.} Ten
patients with two held out per replicate admit 45 distinct holdout pairs, so
every replicate count above that reuses pairs and every interval quoted on the
calibration runs optimistic. Raising the replicate count stabilises Monte Carlo
noise and converges the estimate to the value over those 45 pairs rather than
to a population value. The patient-level uncertainty is unchanged by the
replicate count and is not quantified anywhere in this paper: read the returned
cutpoint as what this procedure gives on this cohort, not as a number with an
interval on it.

\paragraph{One resolution is degenerate by construction.} At the coarsest
resolution every scored cell belongs to the type in question, so the
compositional term sits at exactly zero and only the intrinsic term carries
information. Two intermediate resolutions are the same partition on most
patients, making one measurement reported twice rather than two agreeing
measurements. Measuring the score distribution established both facts rather
than assuming them. Both annotation axes score \emph{immaturity} and the score
gets negated, so ``no marker detected'' becomes the maximum attainable maturity
and the zero-atom sits at the score maximum. Any score that is a function of
the marker counts alone gives those cells the same value, verified across six
formulations. The formulations breaking the tie use information beyond the
markers and order the atom by sequencing depth, which is arm-confounded.

\section{The bound, on labels rare enough for it to bite}
\label{app:rare}

\S\ref{sec:ceiling} reports that no row of our own maturity labels falls in
$0 < p < 1.3516\%$, so the blind spot the bound predicts is unpopulated there.
That is a fact about our four annotation resolutions, not about the data. Both
Lee cohorts ship a published per-cell subtype annotation --- 36 subtypes on SMC,
40 on KUL3 --- and a dozen of those sit under $1.35\%$.

Scoring each subtype as a binary label against sequencing depth, through the
same function, the same QC-filtered cells and the same whole-transcriptome
library size that produced Figure~\ref{fig:ceiling}, gives the table below. Only
the label changes, so the two are directly comparable.

\begin{center}
\small
\begin{tabular}{@{}lrrrr@{}}
\toprule
& \textbf{rows} & \textbf{median $|\rho|$/bound} & \textbf{flagged at 0.20} &
\textbf{above $\tfrac{1}{2}$ bound} \\
\midrule
\multicolumn{5}{@{}l}{\emph{GSE132465 (SMC), 10 paired patients}} \\
\quad $0 < p < 1.3516\%$ & 309 & 0.424 & 0 & 118 \\
\quad $p \geq 1.3516\%$ & 240 & 0.386 & 72 (30\%) & --- \\[2pt]
\multicolumn{5}{@{}l}{\emph{GSE144735 (KUL3), 6 paired patients}} \\
\quad $0 < p < 1.3516\%$ & 261 & 0.502 & 0 & 132 \\
\quad $p \geq 1.3516\%$ & 226 & 0.316 & 49 (22\%) & --- \\[2pt]
\midrule
\multicolumn{5}{@{}l}{\emph{both}} \\
\quad $0 < p < 1.3516\%$ & \textbf{570} & 0.447 & \textbf{0} & \textbf{250} \\
\quad $p \geq 1.3516\%$ & 466 & 0.358 & 121 (26\%) & --- \\
\bottomrule
\end{tabular}
\end{center}

Three things to read off it. The blind rows are \emph{rare}, not degenerate ---
none sits at $p$ exactly $0$ or $1$, which is what separates them from the 58
rows of \S\ref{sec:ceiling}. As a share of what each row could attain they are
no less confounded than the rows the check can see, and on KUL3 they are more so
($0.502$ against $0.316$); the maximum is $0.994$, a label at $99.4\%$ of the
strongest depth association its prevalence permits, reported clean. And the
populations inside the band include the mature enterocyte, goblet and stem-like
compartments this decomposition is about.

Two caveats we do not want read past. These are the authors' annotations, not
our maturity call; that is deliberate, since the claim under test is about the
\emph{statistic} rather than about our labelling, and labels we did not build
are the stronger demonstration --- but it means this shows the check going blind
on cell-type labels generally, not on ours specifically. And published
annotations are themselves derived from expression, so some depth association is
expected. Expected is not the same as reportable: the diagnostic's job was to
report it, and at these prevalences it cannot.
